% BEGIN FINAL TEXT FIX FT02: Abstract
\chapter*{Abstract}
Malicious activity on Ethereum can involve transaction sequences that are
difficult to recognise from individual address summaries. This thesis examines
whether temporal motif fingerprints improve malicious-address detection, how
their usefulness depends on the temporal window, and whether they distinguish
reported behaviour categories. The fingerprints count each address's local
roles in ordered three-transfer candidate patterns. They are evaluated on two
labelled Ethereum datasets using held-out addresses, nine windows from five
minutes to 30 days, and seven orderings of transfers with equal timestamps.

Adding fingerprints to 17 conventional transaction and graph features increases
average precision in all six cohort comparisons on the main dataset, but only
three increases are statistically significant. The independent binary comparison
also shows no significant improvement. Larger windows generally increase both
average precision and fingerprint coverage, without establishing a unique
optimum. Across nine malicious-behaviour categories, fingerprints contain
statistically detectable category information, but classification remains
limited, with macro average precision of 0.1953. Selecting separate windows for
the categories gives no significant improvement over a shared window.

The results show that temporal motif fingerprints describe behavioural
differences, but do not consistently add predictive value beyond conventional
features. The evaluation concerns retrospective address classification, not
direct detection of independently verified coordinated campaigns.
% END FINAL TEXT FIX FT02

